\section{Introduction}

\begin{frame}{Problem Statement}
Consider the usual linear regression problem with noise
\[
b = Ax + \epsilon, \tag{P}\label{P}
\]
where $b \in \R^n$ is the response, $\epsilon \in \R^n$ is the noise, 
$A \in \R^{n \times p}$ is the design matrix, and $x \in \R^p$ is the coefficient vector 
to be estimated.

With the increase in data dimension, sparse learning techniques have gained attention 
for their compacity and interpretability \cite{statsparse,stathigh}.
\end{frame}

\begin{frame}{Sparse Regularized Problem}
Assuming sparsity in the solution (often desired or theoretically expected when $p \gg n$), a natural strategy is to consider the regularized least squares problem
\[
\min_{x \in \R^p} \frac{1}{2} \|Ax-b\|_2^2 + \lambda_0 \|x\|_0 + \lambda_q \|x\|_q^q,\quad q \in \{1,2\}. \tag{R}\label{R}
\]
The $\ell_0$ pseudo-norm $\|x\|_0 = |\{i \mid x_i \neq 0\}|$; $\lambda_0 \ge 0$ controls the trade-off between nonzeros and data fit; $\lambda_q \ge 0$ controls the $\ell_q$ regularization.
\end{frame}

\begin{frame}{Motivation for $\ell_q$ Regularization}
Although $q=1$ also induces sparsity, the additional $\ell_q$ term helps avoid \textit{overfitting} in low Signal-to-Noise Ratio (SNR) regimes \cite{overfitting} via \textit{shrinkage}. 
We denote \eqref{R} with $q=1$ as the $\ell_0\ell_1$ problem and $q=2$ as the $\ell_0\ell_2$ problem.
\end{frame}

\begin{frame}{Approaches to Solve \eqref{R}}
\eqref{R} is NP-hard \cite{nphard}. Three common approaches:
\begin{enumerate}
\item \textbf{Convex/nonconvex proxies}: LASSO \cite{lasso}, MCP \cite{mcp}, SCAD \cite{scad};
\item \textbf{Exact methods}: Mixed Integer Programming (MIP) \cite{mio} – can solve globally for $p \sim 10^7$ when high sparsity is desired, but at high cost;
\item \textbf{Heuristics / approximate algorithms}: stepwise regression, Iterative Hard Thresholding (IHT) \cite{iht}, \texttt{abess} \cite{abess}, greedy/randomized Coordinate Descent (CD) \cite{gcd,rcd}.
\end{enumerate}
\end{frame}

\begin{frame}{The \texttt{L0Learn} Approach}
\cite{fastselect} proposes a method that balances speed and near-optimality: 
\begin{itemize}
\item Seeks \textit{almost optimal} solutions satisfying a local combinatorial exactness property;
\item Small perturbations of the support do not improve the objective;
\item The resulting estimators outperform others in prediction, estimation, and variable selection;
\item The \texttt{L0Learn} implementation is faster than \texttt{ncvreg} \cite{ncvreg} and \texttt{glmnet} \cite{glmnet}.
\end{itemize}
\end{frame}

\begin{frame}{Notation}
\begin{itemize}
\item $T \subset \{1,\dots,p\}$: $U_T \in \R^{p \times |T|}$ is the submatrix of the identity $I_p$ corresponding to indices in $T$;
\item $x_T$: subvector of $x$ with indices $T$;
\item $\Supp(x) \coloneqq \{i \mid x_i \neq 0\}$: support of vector $x$;
\item $\mathcal{C}^{1,1}_L$: functions $g:\R^k \to \R$, $\mathcal{C}^1$ with $L$-Lipschitz gradient;
\item $\lambda_{\max}(B)$: largest eigenvalue of $B$;
\item $[p] \coloneqq \{1,\dots,p\}$.
\end{itemize}
\end{frame}